% Extracted directly from SYS_SAL and SHOTS_SAL, not paraphrased.
\paragraph{SAL keep-biased system instruction.}
\begin{quote}\small
You judge whether the «marked» TARGET sentence is a SELF-ACKNOWLEDGED LIMITATION of the AUTHORS' OWN study, using the SAL annotation criteria. A limitation is a weakness/caveat/bias/constraint the authors admit about THIS trial, fitting a known category: study design; loose or narrow population / convenience sampling; single-center or narrow setting; intervention composite/non-standardized; control (no placebo / active placebo / care-as-usual); outcome measure excluded/imprecise/unvalidated; missing data or drop-out; small sample / underpowered; randomization unbalanced/poor/open; blinding (patient or team not blinded); short duration or follow-up; multiple/post-hoc testing or unadjusted confounding; funding; limited generalizability.

DECISION RULE (keep-biased):

- keep=true if the TARGET states such a limitation, OR states a consequence ('could introduce bias', 'cannot be assessed', 'results may be confounded', 'risk of type 1 error') whose underlying CAUSE (no blinding, small sample, single center, multiple testing, short follow-up...) is evident in the TARGET or its SURROUNDING context. When in doubt and a real weakness of this study is present, keep=true.

- keep=false ONLY for CLEAR non-limitations: restated positive results; clinical recommendations; critique of prior/other work; plain methods description; background; or a bare header like 'This study has several limitations.' with no specific weakness.

\end{quote}
\paragraph{Few-shot pairs.}
\begin{quote}\small
\noindent\textbf{YES}: the lack of a no treatment arm in this study\par
\noindent\textbf{YES}: our study included only a relatively small sample of participants\par
\noindent\textbf{YES}: The absence of masking is most relevant for the assessment of events.\par
\noindent\textbf{YES}: the crossover design and the relatively short treatment periods are limitations\par
\noindent\textbf{YES}: the trial was conducted at a single center\par
\noindent\textbf{YES}: Because outcomes were exploratory, the risk of type 1 error was considerable.\par
\noindent\textbf{YES}: This could have introduced a bias in the results.\par
\noindent\textbf{NO}: This study has several limitations.\par
\noindent\textbf{NO}: The intervention significantly reduced symptoms at 12 weeks.\par
\noindent\textbf{NO}: Clinicians should consider these findings when counseling patients.\par
\noindent\textbf{NO}: Existing methods cannot capture long-range dependencies.\par
\noindent\textbf{NO}: Participants were randomized 1:1 using a computer-generated sequence.\par
\end{quote}
